% P45の方法部品。投稿の中心手法・完成原稿ではない。
% 固定条件と全件監査は docs/experiments/P45_documentation_comparison.md を参照。
\subsection{途中状態を固定した仕様表示の局所比較}
入力の短縮が最終品質と仕事全体の費用に結び付くかを調べるため，
既知の開発用4課題について，第8応答後の同じ会話，環境，Python状態から
全表示と仕様返却の選択表示を各1回継続する．
課題を成功率や削減量に応じて入れ替えず，全8条件を集計する．
選択表示では，保存された表示式，実際のAPI要求・返却，出力全体が
限定した構文規則で一致した古い仕様返却部分だけを原文参照へ置き換える．
作業データ，見出し，操作，最新出力を残し，
その後の実行結果は通常どおり追記する．
この比較は単一作業役の通常の選択省略であり，複数担当の通信削減を測らない．

両条件で同じ残予算通知と残量に応じた出力上限制御を用いる．
継続の上限は16モデル呼出し，入力と出力の合計120,000トークン，480秒とし，
API上限は共通前半を含む元課題の上限を維持する．
各呼出しでは全入力を計数し，残トークンから入力を引いた値と
4,096の小さい方を出力枠とする．
最初の生成は両条件とも推論キャッシュを再利用せず，実測値で確認する．
以後の同条件内での再利用は認める．
初回キャッシュ条件の整合は，生成の完全な決定性を保証するものではない．

継続費用は通知，表示構築，入力計数，生成，ツール操作を含み，
失敗した条件も除外しない．
共通前半，実験用の状態復元，モデルロードを別に記録する．
前半を各条件に1回加えた論理的な全仕事の費用を併記するが，
新たな端から端までの実行を測定した値とは区別する．
使用量やAPI作用が不明，復元不一致，実行・評価の故障が発生した場合は
予定全体を停止し，不明量と未実行条件を保存する．

最終成否は生成側へ返さない独立評価器で判定する．
全8条件の評価，既知の使用量，原記録監査の合格を前提に，
全表示で成功した課題の喪失がなく，継続総トークンが少ないことを
次の設計へ進む開発上の基準とする．
同じ4既知課題の各1回であり，この基準を満たしても
母集団における統計的非劣性や方式の新規性を証明したとはしない．
